\section{Discussion and Conclusion}

Our results show that jointly preserving within-dataset affinity structure while aligning cells across batches leads to improved metacell construction. The combination of a community-preserving objective with reconstruction-based alignment produces metacells that better reflect biological community structure, enabling more faithful aggregation across datasets. Importantly, reconstruction-only objectives are insufficient to preserve affinity-defined structure: due to variations in manifold density, rare cell populations can be absorbed into denser regions of the latent space, and spatially meaningful relationships may be lost. Incorporating an affinity-guided objective mitigates these issues by explicitly enforcing community structure.

These findings highlight limitations in existing approaches. Most methods operate within individual datasets, restricting the number of samples available per cellular state and limiting coverage of rare populations. Reconstruction-based models align cells across batches but do not preserve community structure, while affinity-based methods preserve structure but remain restricted to within-batch aggregation. scProto addresses both by combining affinity guidance with cross-batch alignment.

Despite these advantages, the number of metacells $K$ must be specified in advance, which may not be optimal across datasets. An important direction for future work is to develop methods that automatically determine $K$ and jointly learn similarity structure and representation.

The framework operates on an affinity graph constructed in PCA space, which reflects both biological variation and technical batch effects—a deliberate design choice, since the batch-conditional encoder is responsible for disentangling these signals during training.
